\section{Related Work, Limitations, and Conclusion}
\label{sec:discussion}
\label{sec:conclusion}
\paragraph{Proof auditing and error localization.}
Step-level proof benchmarks such as ProcessBench~\citep{zheng2025processbench}, PRMBench~\citep{song2025prmbench}, and Hard2Verify~\citep{pandit2025hard2verify} evaluate step-level correctness or error localization. Closer to formal verification, ProofFlow tracks step dependencies during autoformalization \citep{cabral2026proofflow}, Critique-Resilient Benchmarking evaluates adversarial exchanges through checked objections \citep{marro2026crb}, and FaithSieve decomposes informal proofs into local claims that combine Lean checks with semantic alignment assessments \citep{wang2026faithsieve}. Co-PEER complements these evaluations by storing the concrete counterexample, its Lean certificate, and the human or model review of source correspondence in a persistent dossier that backs each benchmark reference line.

\paragraph{Counterexample search and formal translation.}
CounterMath evaluates conceptual reasoning through counterexamples~\citep{li2025countermath}, while Learning to Disprove~\citep{li2026disprove} and Lean Disprove~\citep{ondras2026disprove} evaluate automated counterexample search. FormalAlign and recent audits of formal math benchmarks show that a statement checked by Lean can still misrepresent the informal problem it translates \citep{lu2025formalalign,ammanamanchi2026faults}. Co-PEER brings those two threads together by separating Lean's check of the formal counterexample from the review of source correspondence. In the spirit of \citet{lakatos1976proofs}, it also uses checked counterexamples to repair proofs and characterize families of failures (see Appendix~\ref{sec:related} for a comparison table of related resources).

\paragraph{Limitations and open questions.}
Each part of our evaluation has a defined scope. First, BugBank stores successful objections, and the localization benchmark evaluates texts known to contain an error. Measuring end-to-end auditing reliability requires correct-proof controls for false accusations, a defined starting population to estimate how often search, formalization, and source review succeed, and matched computation budgets on unseen sources to compare efficiency. Second, human reviewers and model graders can share blind spots, and annotated reference lines need not exhaust every valid objection in a proof. Third, lower output-token counts under counterexample-first search do not establish lower total computational or human cost once candidate selection, failed searches, Lean formalization, source interpretation, and human review are included. Fourth, the repair and retrieval studies use small task sets and show model-dependent outcomes, so broader trials are needed to determine when supplied witnesses or retrieved dossiers improve an argument.

\paragraph{Conclusion.}
A counterexample is most useful when a reader can verify both that the witness breaks the formalized claim and that the formalized claim matches the source argument. Co-PEER captures both checks, together with scope and revision history, in an objection dossier, and BugBank's human-approved core of 147 dossiers supports our 100-text localization benchmark. Separate studies test search instructions, repair tasks, and retrieval of earlier objections. Beyond auditing single steps, analyzing why a counterexample works can turn an isolated failure into a complete correction or a theorem for an infinite family.
